\chapter{The G-Machine}
\vspace{2cm}
%\vspace{1cm}

The heart of any graph reducer is the implementation of function application.
In Chapter~12 we saw that a lambda abstraction can be applied to an argument
by constructing an instance of the body of the abstraction with substitutions
made for occurrences of the formal parameter. Unfortunately, this involved
an inefficient traversal of the tree representing the body of the abstraction,
and the presence of free variables seemed to make a more efficient implementation rather difficult.

With this in mind, we developed the supercombinator transformation in
Chapter~13, which yielded particularly simple lambda abstractions (the
supercombinators), which had no free variables. This simplified the process of
instantiating the body of such an abstraction, but at the (minor) price of
having to substitute for several variables at once. However, the principal
incentive for developing the supercombinator transformation was the hope of
compiling the body of a supercombinator to a fixed sequence of instructions
which, when executed, would construct an instance of its body.

The payoff comes in this chapter, in which we will examine the G-machine,
an extremely fast implementation of graph reduction based on
supercombinator compilation. The G-machine was developed at the Chalmers
Institute of Technology, G\"oteborg, Sweden, by Johnsson and Augustsson.
This chapter and the subsequent three chapters draw heavily on the
G-machine papers [Johnsson, 1984; Augustsson, 1984]. Many of the ideas in
these chapters are theirs, and not all of them have appeared in the published
literature.

The development of the G-machine is presented informally, but it would be
an interesting exercise to give a formal proof of its correctness [Lester, 1985].

\section{Using an Intermediate Code}

Once we have decided to compile supercombinator bodies to a sequence of
instructions we have to decide on the language in which the instructions
should be written. It would be possible to produce, say, VAX machine code
directly, but this approach suffers from two disadvantages. Firstly, we would
have to start all over again if we want to generate code for some other
machine, and, secondly, we would be in danger of mixing up the issues of how
to compile supercombinators to a sequential code with issues of how best to
exploit particular features of the VAX.

This is not a new problem, and a common solution is to define an
\textit{intermediate code}, which can be regarded as the machine code for an abstract
sequential machine. Then the compilation process can be split into two parts:
first generate the intermediate code, and then generate target code for a
particular machine from the intermediate code. Changing the code generator
to generate code for a different target machine is then relatively easy, and
improvements made in the compilation to intermediate code benefit all such
code generators. Examples of this approach include Pascal's P-code [Clark,
1981], BCPL's O-code [Richards, 1971] and Portable Standard Lisp's
C-macros [Griss and Hearn, 1981].

\subsection{G-code and the G-machine Compiler}

For these reasons, the designers of the G-machine defined an intermediate
code called G-code, into which supercombinator bodies are compiled. The
compiler for the G-machine follows a sequence similar to that described in the
first two parts of this book. In particular:

\begin{numbered}
    \item The source language is a variant of ML with lazy evaluation semantics,
    called Lazy ML (or LML).
    \item Early phases of the compiler perform type-checking, compile pattern-matching and do dependency analysis. At this stage the program has
    been translated to the lambda calculus (augmented with \ml{let} and \ml{letrec}).
    \item A lambda-lifter transforms the program to supercombinator form. The
    full laziness optimization is not performed, but this feature could easily be
    added.
    \item Now the supercombinators are compiled to G-code.
    \item Finally, machine code for the target machine is generated from the
    G-code.
\end{numbered}

\noindent
Figure~18.1 shows the structure of the G-machine compiler.

Our description of the G-machine compiler falls into three parts:

\begin{numbered}
    \item a description of the compilation algorithm which translates the source
    language into the intermediate code;
    \item a description of the intermediate code itself, giving a precise description
    of what each instruction does;
    \item a description of the code generator.
\end{numbered}


\boxedfigure{
    \centering
    \begin{tabular}{@{}c@{\hspace{1.5cm}}l@{}}
        Source program & \\
        $\downarrow$ & \\
        \fbox{Compile to lambda expressions} & \\
        $\downarrow$ & \\
        Lambda calculus plus \ml{let} and \ml{letrec} & \\
        $\downarrow$ & \\
        \fbox{Lambda-lifting} & \\
        $\downarrow$ & \\
        Set of supercombinator definitions & \\
        $\downarrow$ & \\
        \fbox{Compilation to G-code} & Chapters 18, 20, 21 \\
        $\downarrow$ & \\
        G-code program & Chapters 18, 19 \\
        $\downarrow$ & \\
        \fbox{Code generation} & Chapter 19 \\
        $\downarrow$ & \\
        Target machine code & \\
    \end{tabular}
}{Structure of the G-machine compiler}

\noindent
We will discuss the first of these parts in this chapter and the latter two in the
next chapter. First, however, we will mention some related work.

\subsection{Other Fast Sequential Implementations of Lazy Languages}

The implementation of Ponder [Fairbairn, 1982], developed by Fairbairn and
Wray, is based on a similar approach to the G-machine. The Ponder Abstract
Machine (PAM) is at least as sophisticated as the G-machine, though they
were developed independently, and is described in Fairbairn's thesis
[Fairbairn, 1985; Fairbairn and Wray, 1986]. An interesting development of
this work is a cross-compiler which compiles Ponder abstract machine
instructions into SKIM microcode [Elworthy, 1985].

A related approach, though one which diverges from graph reduction, is to
use a lexically scoped dialect of Lisp, such as Scheme [Steele and Sussman,
1978] or T [Rees and Adams, 1982], as an intermediate code. This takes
advantage of the immense amount of effort which has been spent on building
fast Lisp implementations, and is the approach taken by Hudak [Hudak and
Kranz, 1984].

A fast VAX implementation of Hope [Burstall \textit{et al.}, 1980] based on an
intermediate code called FP/M has recently been developed at Imperial
College [Field, 1985] (remember, however, that Hope is a strict language).

\section{An Example of G-machine Execution}

We begin with an example, to give the flavor of the G-machine. Consider the
Miranda program

\combinatorbox{
    \ml{from n = n : from (succ n)} \\
    \indent\ml{succ n = n+1}
}{
    \ml{from (succ 0)}
}

\noindent
It generates the infinite list \ml{[1,2,3,\ldots,]}. The functions \ml{from} and \ml{succ} are
supercombinators already, so the lambda-lifting is trivial, yielding

\combinatorbox{
    \ml{\$from n = CONS n (\$from (\$succ n))} \\
    \indent\ml{\$succ n = + n 1} \\
    \indent\ml{\$Prog = \$from (\$succ 0)}
}{
    \ml{\$Prog}
}

\noindent
The G-machine uses a stack, and execution begins with a pointer to the initial
graph on top of the stack (Figure~18.2(a)). The spine is then unwound, exactly
as previously discussed in Section~11.6, without using pointer-reversal. The
difference comes when the spine has been completely unwound, so that there
is a pointer to \ml{\$from} on the stack (see Figure~18.2(b)). By following this
pointer the machine extracts

\begin{numbered}
    \item the number of arguments expected by \ml{\$from} (one in this case);
    \item the starting address for the code for \ml{\$from}.
\end{numbered}

\noindent
First it checks that there are enough arguments on the stack for \ml{\$from} to
execute, and finds that there are. It then rearranges the top of the stack
slightly (see the transition from (b) to (c) in Figure~18.2) and then jumps to the
code for \ml{\$from}. The rearrangement of the top of the stack puts a pointer to the
argument to \ml{\$from} on top of the stack. We will discuss the stack rearrangement in more detail later. Notice also that the machine \textit{jumps} to \ml{\$from} rather
than \textit{calling} it. An instruction at the end of \ml{\$from} will complete evaluation of
the graph after the \ml{\$from} reduction is done.


\begin{figure}[h]
    \centering

    {%
        \setlength{\fboxrule}{1pt}%
        \setlength{\fboxsep}{0pt}%
        \fbox{%
%            \begin{minipage}{0.9\textwidth}
                \sffamily\footnotesize
%                \setlength{\parindent}{10pt}
                \setlength{\parskip}{0mm plus 0mm minus 0mm}
                \centering\noindent

\begin{tabular}{ccc}

    % a
    \begin{tikzpicture}[
        baseline=(current bounding box.north),
        font=\sffamily\footnotesize,
        cell/.style={draw, minimum width=1.3cm, minimum height=0.5cm, inner sep=0pt},
        bullet/.style={circle, fill=black, inner sep=1.5pt},
        ptr/.style={-{Stealth[length=2.5mm,width=1.5mm]}, thick},
        apply/.style={inner sep=1pt},
        lbl/.style={inner sep=1pt}
        ]

        % Stack: a single cell with a bullet at its center
        \node[cell] (c1) at (0,0) {};
        \node[bullet] (s1) at (c1.center) {};

        % Tree (hand-placed nodes); root @ is at the same height as the bullet
        \node[apply] (ap1)  at (2.4,0)    {\ml{@}};
        \node[lbl]   (from) at (1.7,-0.7) {\ml{\$from}};
        \node[apply] (ap2)  at (2.7,-0.7) {\ml{@}};
        \node[lbl]   (succ) at (2.2,-1.5) {\ml{\$succ}};
        \node[lbl]   (zero) at (3.1,-1.5) {\ml{0}};

        % Tree edges
        \draw (ap1) -- (from);
        \draw (ap1) -- (ap2);
        \draw (ap2) -- (succ);
        \draw (ap2) -- (zero);

        % Pointer from the stack bullet to the root @ node: one straight segment
        \draw[ptr] (s1) -- (ap1.west);
    \end{tikzpicture}
    &

    % b
    \begin{tikzpicture}[
        baseline=(current bounding box.north),
        font=\sffamily\footnotesize,
        cell/.style={draw, minimum width=1.3cm, minimum height=0.5cm, inner sep=0pt},
        bullet/.style={circle, fill=black, inner sep=1.5pt},
        ptr/.style={-{Stealth[length=2.5mm,width=1.5mm]}, thick},
        apply/.style={inner sep=1pt},
        lbl/.style={inner sep=1pt}
        ]

        % Stack: two cells, each with a bullet at its center
        \node[cell] (c1) at (0,0)    {};
        \node[cell] (c2) at (0,-0.5) {};
        \node[bullet] (s1) at (c1.center) {};
        \node[bullet] (s2) at (c2.center) {};

        % Tree (hand-placed nodes); root @ is at the same height as the first bullet
        \node[apply] (ap1)  at (2.4,0)    {\ml{@}};
        \node[lbl]   (from) at (1.7,-0.7) {\ml{\$from}};
        \node[apply] (ap2)  at (2.7,-0.7) {\ml{@}};
        \node[lbl]   (succ) at (2.2,-1.5) {\ml{\$succ}};
        \node[lbl]   (zero) at (3.1,-1.5) {\ml{0}};

        % Tree edges
        \draw (ap1) -- (from);
        \draw (ap1) -- (ap2);
        \draw (ap2) -- (succ);
        \draw (ap2) -- (zero);

        % Pointer from the first bullet to the root @ node
        \draw[ptr] (s1) -- (ap1.west);

        % Pointer from the second bullet to $from (curved, as in the scan)
        \draw[ptr] (s2) to[out=0, in=180] (from.west);
    \end{tikzpicture}

    &

    % c
    \begin{tikzpicture}[
        baseline=(current bounding box.north),
        font=\sffamily\footnotesize,
        cell/.style={draw, minimum width=1.3cm, minimum height=0.5cm, inner sep=0pt},
        bullet/.style={circle, fill=black, inner sep=1.5pt},
        ptr/.style={-{Stealth[length=2.5mm,width=1.5mm]}, thick},
        apply/.style={inner sep=1pt},
        lbl/.style={inner sep=1pt}
        ]

        % Stack: two cells, each with a bullet at its center
        \node[cell] (c1) at (0,0)    {};
        \node[cell] (c2) at (0,-0.5) {};
        \node[bullet] (s1) at (c1.center) {};
        \node[bullet] (s2) at (c2.center) {};

        % Tree (hand-placed nodes)
        \node[apply] (ap1)  at (2.4,0)    {\ml{@}};
        \node[lbl]   (from) at (1.7,-0.7) {\ml{\$from}};
        \node[apply] (ap2)  at (2.7,-0.7) {\ml{@}};
        \node[lbl]   (succ) at (2.2,-1.5) {\ml{\$succ}};
        \node[lbl]   (zero) at (3.1,-1.5) {\ml{0}};

        % Tree edges
        \draw (ap1) -- (from);
        \draw (ap1) -- (ap2);
        \draw (ap2) -- (succ);
        \draw (ap2) -- (zero);

        % Pointer from the first bullet to the root @ node
        \draw[ptr] (s1) -- (ap1.west);

        % Pointer from the second bullet to the lower @ node: curves down and
        % under $from, then rises into ap2 from the lower left
        \draw[ptr] (s2) to[out=-60, in=-150] (ap2.south west);
    \end{tikzpicture}

    \\[2.5cm]
    (a) &
    (b) \ml{UNWIND} &
    (c) Rearrange stack \\[.5cm]

    % d

    \begin{tikzpicture}[
        baseline=(current bounding box.north),
        font=\sffamily\footnotesize,
        cell/.style={draw, minimum width=1.3cm, minimum height=0.5cm, inner sep=0pt},
        bullet/.style={circle, fill=black, inner sep=1.5pt},
        ptr/.style={-{Stealth[length=2.5mm,width=1.5mm]}, thick},
        apply/.style={inner sep=1pt},
        lbl/.style={inner sep=1pt}
        ]

        % Stack: three cells, each with a bullet at its center
        \node[cell] (c1) at (0,0)    {};
        \node[cell] (c2) at (0,-0.5) {};
        \node[cell] (c3) at (0,-1.0) {};
        \node[bullet] (s1) at (c1.center) {};
        \node[bullet] (s2) at (c2.center) {};
        \node[bullet] (s3) at (c3.center) {};

        % Tree (hand-placed nodes)
        \node[apply] (ap1)  at (2.4,0)    {\ml{@}};
        \node[lbl]   (from) at (1.7,-0.7) {\ml{\$from}};
        \node[apply] (ap2)  at (2.7,-0.7) {\ml{@}};
        \node[lbl]   (succ) at (2.2,-1.5) {\ml{\$succ}};
        \node[lbl]   (zero) at (3.1,-1.5) {\ml{0}};

        % Tree edges
        \draw (ap1) -- (from);
        \draw (ap1) -- (ap2);
        \draw (ap2) -- (succ);
        \draw (ap2) -- (zero);

        % Pointer from the first bullet to the root @ node
        \draw[ptr] (s1) -- (ap1.west);

        % Second pointer: shallow curve into the west side of ap2
        \draw[ptr] (s2) to[out=-30, in=-160] (ap2.west);

        % Third pointer: slightly deeper curve into the south-west side of ap2
        \draw[ptr] (s3) to[out=-35, in=-140] (ap2.south west);
    \end{tikzpicture}


    &

    % e
    \begin{tikzpicture}[
        baseline=(current bounding box.north),
        font=\sffamily\footnotesize,
        cell/.style={draw, minimum width=1.3cm, minimum height=0.5cm, inner sep=0pt},
        bullet/.style={circle, fill=black, inner sep=1.5pt},
        ptr/.style={-{Stealth[length=2.5mm,width=1.5mm]}, thick},
        apply/.style={inner sep=1pt},
        lbl/.style={inner sep=1pt}
        ]

        % Stack: four cells, each with a bullet at its center
        \node[cell] (c1) at (0,0)    {};
        \node[cell] (c2) at (0,-0.5) {};
        \node[cell] (c3) at (0,-1.0) {};
        \node[cell] (c4) at (0,-1.5) {};
        \node[bullet] (s1) at (c1.center) {};
        \node[bullet] (s2) at (c2.center) {};
        \node[bullet] (s3) at (c3.center) {};
        \node[bullet] (s4) at (c4.center) {};

        % Tree (hand-placed nodes)
        \node[apply] (ap1)  at (2.4,0)    {\ml{@}};
        \node[lbl]   (from) at (1.7,-0.7) {\ml{\$from}};
        \node[apply] (ap2)  at (2.7,-0.7) {\ml{@}};
        \node[lbl]   (succ) at (2.2,-1.5) {\ml{\$succ}};
        \node[lbl]   (zero) at (3.1,-1.5) {\ml{0}};

        % Second tree: a single node beneath the first
        \node[lbl]   (succ2) at (1.9,-2.3) {\ml{\$succ}};

        % Tree edges
        \draw (ap1) -- (from);
        \draw (ap1) -- (ap2);
        \draw (ap2) -- (succ);
        \draw (ap2) -- (zero);

        % Pointer from the first bullet to the root @ node
        \draw[ptr] (s1) -- (ap1.west);

        % Second pointer: shallow curve into the west side of ap2
        \draw[ptr] (s2) to[out=-30, in=-160] (ap2.west);

        % Third pointer: slightly deeper curve into the south-west side of ap2
        \draw[ptr] (s3) to[out=-15, in=-140] (ap2.south west);

        % Fourth pointer: to the new $succ node
        \draw[ptr] (s4) to[out=0, in=120] (succ2.north);
    \end{tikzpicture}

    &

    % f
    \begin{tikzpicture}[
        baseline=(current bounding box.north),
        font=\sffamily\footnotesize,
        cell/.style={draw, minimum width=1.3cm, minimum height=0.5cm, inner sep=0pt},
        bullet/.style={circle, fill=black, inner sep=1.5pt},
        ptr/.style={-{Stealth[length=2.5mm,width=1.5mm]}, thick},
        apply/.style={inner sep=1pt},
        lbl/.style={inner sep=1pt}
        ]

        % Stack: three cells, each with a bullet at its center
        \node[cell] (c1) at (0,0)    {};
        \node[cell] (c2) at (0,-0.5) {};
        \node[cell] (c3) at (0,-1.0) {};
        \node[bullet] (s1) at (c1.center) {};
        \node[bullet] (s2) at (c2.center) {};
        \node[bullet] (s3) at (c3.center) {};

        % First tree (hand-placed nodes)
        \node[apply] (ap1)  at (2.4,0)    {\ml{@}};
        \node[lbl]   (from) at (1.7,-0.7) {\ml{\$from}};
        \node[apply] (ap2)  at (2.7,-0.7) {\ml{@}};
        \node[lbl]   (succ) at (2.2,-1.5) {\ml{\$succ}};
        \node[lbl]   (zero) at (3.1,-1.5) {\ml{0}};

        % Second tree: an apply node over a single $succ leaf
        \node[apply] (ap3)   at (1.3,-1.8) {\ml{@}};
        \node[lbl]   (succ2) at (0.9,-2.4) {\ml{\$succ}};

        % First tree edges
        \draw (ap1) -- (from);
        \draw (ap1) -- (ap2);
        \draw (ap2) -- (succ);
        \draw (ap2) -- (zero);

        % Second tree edge
        \draw (ap3) -- (succ2);

        % Pointer from the first bullet to the root @ node
        \draw[ptr] (s1) -- (ap1.west);

        % Second pointer: shallow curve into the west side of ap2
        \draw[ptr] (s2) to[out=-30, in=-160] (ap2.west);

        % Third pointer: now ends at the new apply node ap3
        \draw[ptr] (s3) to[out=-15, in=120] (ap3.north);

        % Fourth pointer: from the right of ap3, curling up and into ap2
        \draw[ptr] (ap3.east) to[out=-30, in=-140] (ap2.south west);
    \end{tikzpicture}

    \\[3.4cm]
    (d) \ml{PUSH 0} &
    (e) \ml{PUSHGLOBAL \$succ} &
    (f) \ml{MKAP} \\[.5cm]

    %g

    \begin{tikzpicture}[
        baseline=(current bounding box.north),
        font=\sffamily\footnotesize,
        cell/.style={draw, minimum width=1.3cm, minimum height=0.5cm, inner sep=0pt},
        bullet/.style={circle, fill=black, inner sep=1.5pt},
        ptr/.style={-{Stealth[length=2.5mm,width=1.5mm]}, thick},
        apply/.style={inner sep=1pt},
        lbl/.style={inner sep=1pt}
        ]

        % Stack: four cells, each with a bullet at its center
        \node[cell] (c1) at (0,0)    {};
        \node[cell] (c2) at (0,-0.5) {};
        \node[cell] (c3) at (0,-1.0) {};
        \node[cell] (c4) at (0,-1.5) {};
        \node[bullet] (s1) at (c1.center) {};
        \node[bullet] (s2) at (c2.center) {};
        \node[bullet] (s3) at (c3.center) {};
        \node[bullet] (s4) at (c4.center) {};

        % First tree (hand-placed nodes)
        \node[apply] (ap1)  at (2.4,0)    {\ml{@}};
        \node[lbl]   (from) at (1.7,-0.7) {\ml{\$from}};
        \node[apply] (ap2)  at (2.7,-0.7) {\ml{@}};
        \node[lbl]   (succ) at (2.2,-1.5) {\ml{\$succ}};
        \node[lbl]   (zero) at (3.1,-1.5) {\ml{0}};

        % Second tree: an apply node over a single $succ leaf
        \node[apply] (ap3)   at (1.3,-1.8) {\ml{@}};
        \node[lbl]   (succ2) at (0.9,-2.4) {\ml{\$succ}};

        % Third tree: a single $from node directly below the fourth cell
        \node[lbl]   (from2) at (0,-2.8) {\ml{\$from}};

        % First tree edges
        \draw (ap1) -- (from);
        \draw (ap1) -- (ap2);
        \draw (ap2) -- (succ);
        \draw (ap2) -- (zero);

        % Second tree edge
        \draw (ap3) -- (succ2);

        % Pointer from the first bullet to the root @ node
        \draw[ptr] (s1) -- (ap1.west);

        % Second pointer: shallow curve into the west side of ap2
        \draw[ptr] (s2) to[out=-30, in=-160] (ap2.west);

        % Third pointer: to the apply node ap3
        \draw[ptr] (s3) to[out=-15, in=120] (ap3.north);

        % Fourth pointer: from the right of ap3, curling up and into ap2
        \draw[ptr] (ap3.east) to[out=-30, in=-140] (ap2.south west);

        % Fifth pointer: straight down from the new cell's bullet to $from
        \draw[ptr] (s4) -- (from2.north);
    \end{tikzpicture}

    &

    %h

    \begin{tikzpicture}[
        baseline=(current bounding box.north),
        font=\sffamily\footnotesize,
        cell/.style={draw, minimum width=1.3cm, minimum height=0.5cm, inner sep=0pt},
        bullet/.style={circle, fill=black, inner sep=1.5pt},
        ptr/.style={-{Stealth[length=2.5mm,width=1.5mm]}, thick},
        apply/.style={inner sep=1pt},
        lbl/.style={inner sep=1pt}
        ]

        % Stack: three cells, each with a bullet at its center
        \node[cell] (c1) at (0,0)    {};
        \node[cell] (c2) at (0,-0.5) {};
        \node[cell] (c3) at (0,-1.0) {};
        \node[bullet] (s1) at (c1.center) {};
        \node[bullet] (s2) at (c2.center) {};
        \node[bullet] (s3) at (c3.center) {};

        % First tree (hand-placed nodes)
        \node[apply] (ap1)  at (2.4,0)    {\ml{@}};
        \node[lbl]   (from) at (1.7,-0.7) {\ml{\$from}};
        \node[apply] (ap2)  at (2.7,-0.7) {\ml{@}};
        \node[lbl]   (succ) at (2.2,-1.5) {\ml{\$succ}};
        \node[lbl]   (zero) at (3.1,-1.5) {\ml{0}};

        % Second tree: ($from ($succ)) combined into one tree
        \node[apply] (ap3)   at (1.0,-1.8) {\ml{@}};
        \node[lbl]   (from2) at (0.5,-2.4) {\ml{\$from}};
        \node[apply] (ap4)   at (1.5,-2.4) {\ml{@}};
        \node[lbl]   (succ2) at (1.2,-3.1) {\ml{\$succ}};

        % First tree edges
        \draw (ap1) -- (from);
        \draw (ap1) -- (ap2);
        \draw (ap2) -- (succ);
        \draw (ap2) -- (zero);

        % Second tree edges
        \draw (ap3) -- (from2);
        \draw (ap3) -- (ap4);
        \draw (ap4) -- (succ2);

        % Pointer from the first bullet to the root @ node
        \draw[ptr] (s1) -- (ap1.west);

        % Second pointer: shallow curve into the west side of ap2
        \draw[ptr] (s2) to[out=-30, in=-160] (ap2.west);

        % Third pointer: to the root of the new tree
        \draw[ptr] (s3) to[out=-15, in=120] (ap3.north);

        % Pointer from the new lower apply node, looping right and up into ap2
        \draw[ptr] (ap4.east) to[out=-40, in=-140] (ap2.south west);
    \end{tikzpicture}

    &
    % i
    \begin{tikzpicture}[
        baseline=(current bounding box.north),
        font=\sffamily\footnotesize,
        cell/.style={draw, minimum width=1.3cm, minimum height=0.5cm, inner sep=0pt},
        bullet/.style={circle, fill=black, inner sep=1.5pt},
        ptr/.style={-{Stealth[length=2.5mm,width=1.5mm]}, thick},
        apply/.style={inner sep=1pt},
        lbl/.style={inner sep=1pt}
        ]

        % Stack: four cells, each with a bullet at its center
        \node[cell] (c1) at (0,0)    {};
        \node[cell] (c2) at (0,-0.5) {};
        \node[cell] (c3) at (0,-1.0) {};
        \node[cell] (c4) at (0,-1.5) {};
        \node[bullet] (s1) at (c1.center) {};
        \node[bullet] (s2) at (c2.center) {};
        \node[bullet] (s3) at (c3.center) {};
        \node[bullet] (s4) at (c4.center) {};

        % First tree (hand-placed nodes)
        \node[apply] (ap1)  at (2.4,0)    {\ml{@}};
        \node[lbl]   (from) at (1.7,-0.7) {\ml{\$from}};
        \node[apply] (ap2)  at (2.7,-0.7) {\ml{@}};
        \node[lbl]   (succ) at (2.2,-1.5) {\ml{\$succ}};
        \node[lbl]   (zero) at (3.1,-1.5) {\ml{0}};

        % Second tree: ($from ($succ)) combined into one tree
        \node[apply] (ap3)   at (1.0,-1.8) {\ml{@}};
        \node[lbl]   (from2) at (0.5,-2.4) {\ml{\$from}};
        \node[apply] (ap4)   at (1.5,-2.4) {\ml{@}};
        \node[lbl]   (succ2) at (1.2,-3.1) {\ml{\$succ}};

        % First tree edges
        \draw (ap1) -- (from);
        \draw (ap1) -- (ap2);
        \draw (ap2) -- (succ);
        \draw (ap2) -- (zero);

        % Second tree edges
        \draw (ap3) -- (from2);
        \draw (ap3) -- (ap4);
        \draw (ap4) -- (succ2);

        % Pointer from the first bullet to the root @ node
        \draw[ptr] (s1) -- (ap1.west);

        % Second pointer: shallow curve into the west side of ap2
        \draw[ptr] (s2) to[out=-30, in=-160] (ap2.west);

        % Third pointer: to the root of the new tree
        \draw[ptr] (s3) to[out=-15, in=120] (ap3.north);

        % Fourth pointer: from the new cell, a shallow curve into ap2 again
        \draw[ptr] (s4) to[out=20, in=-150] (ap2.south west);

        % Pointer from the new lower apply node, looping right and up into ap2
        \draw[ptr] (ap4.east) to[out=-40, in=-140] (ap2.south west);
    \end{tikzpicture}

    \\[3.7cm]

    (g) \ml{PUSHGLOBAL \$from} &
    (h) \ml{MKAP} &
    (i) \ml{PUSH 1} \\[.5cm]

    % j


    \begin{tikzpicture}[
        baseline=(current bounding box.north),
        font=\sffamily\footnotesize,
        cell/.style={draw, minimum width=1.3cm, minimum height=0.5cm, inner sep=0pt},
        bullet/.style={circle, fill=black, inner sep=1.5pt},
        ptr/.style={-{Stealth[length=2.5mm,width=1.5mm]}, thick},
        apply/.style={inner sep=1pt},
        lbl/.style={inner sep=1pt}
        ]

        % Stack: three cells, each with a bullet at its center
        \node[cell] (c1) at (0,0)    {};
        \node[cell] (c2) at (0,-0.5) {};
        \node[cell] (c3) at (0,-1.0) {};
        \node[bullet] (s1) at (c1.center) {};
        \node[bullet] (s2) at (c2.center) {};
        \node[bullet] (s3) at (c3.center) {};

        % First tree (hand-placed nodes)
        \node[apply] (ap1)  at (2.4,0)    {\ml{@}};
        \node[lbl]   (from) at (1.7,-0.7) {\ml{\$from}};
        \node[apply] (ap2)  at (2.7,-0.7) {\ml{@}};
        \node[lbl]   (succ) at (2.2,-1.5) {\ml{\$succ}};
        \node[lbl]   (zero) at (3.1,-1.5) {\ml{0}};

        % Second tree
        \node[apply] (ap3)   at (1.0,-2.4) {\ml{@}};
        \node[lbl]   (from2) at (0.5,-3.0) {\ml{\$from}};
        \node[apply] (ap4)   at (1.5,-3.0) {\ml{@}};
        \node[lbl]   (succ2) at (1.2,-3.7) {\ml{\$succ}};

        % CONS root: stacked double bullet, up and to the left of ap3
        \node[bullet] (cons1) at (0.5,-1.75) {};
        \node[bullet] (cons2) at (0.5,-1.95) {};

        % First tree edges
        \draw (ap1) -- (from);
        \draw (ap1) -- (ap2);
        \draw (ap2) -- (succ);
        \draw (ap2) -- (zero);

        % Second tree edges; the CONS edge branches down and to the right
        \draw (cons2) -- (ap3);
        \draw (ap3) -- (from2);
        \draw (ap3) -- (ap4);
        \draw (ap4) -- (succ2);

        % Pointer from the first bullet to the root @ node
        \draw[ptr] (s1) -- (ap1.west);

        % Second pointer: shallow curve into the west side of ap2
        \draw[ptr] (s2) to[out=-30, in=-160] (ap2.west);

        % Third pointer: to the top of the CONS root
        \draw[ptr] (s3) to[out=-15, in=120] (cons1.north);

        % Pointer from the CONS root into ap2
        \draw[ptr] (cons1.east) to[out=-10, in=-150] (ap2.south west);

        % Pointer from the lower apply node, looping right and up into ap2
        \draw[ptr] (ap4.east) to[out=-40, in=-140] (ap2.south west);
    \end{tikzpicture}

     &
    % k
    \begin{tikzpicture}[
        baseline=(current bounding box.north),
        font=\sffamily\footnotesize,
        cell/.style={draw, minimum width=1.3cm, minimum height=0.5cm, inner sep=0pt},
        bullet/.style={circle, fill=black, inner sep=1.5pt},
        ptr/.style={-{Stealth[length=2.5mm,width=1.5mm]}, thick},
        apply/.style={inner sep=1pt},
        lbl/.style={inner sep=1pt}
        ]

        % Stack: two cells, each with a bullet at its center
        \node[cell] (c1) at (0,0)    {};
        \node[cell] (c2) at (0,-0.5) {};
        \node[bullet] (s1) at (c1.center) {};
        \node[bullet] (s2) at (c2.center) {};

        % CONS root: stacked double bullet
        \node[bullet] (cons1) at (1.7,0.1)  {};
        \node[bullet] (cons2) at (1.7,-0.1) {};

        % Tree on a regular grid: right child +0.4 (or +0.45) in x, left leaf -0.5 in x, 0.7 down per level
        \node[apply] (apa)   at (2.3,-0.6)  {\ml{@}};
        \node[lbl]   (from)  at (1.8,-1.3)  {\ml{\$from}};
        \node[apply] (apb)   at (2.7,-1.3)  {\ml{@}};
        \node[lbl]   (succ1) at (2.2,-2.0)  {\ml{\$succ}};
        \node[apply] (apc)   at (3.1,-2.0)  {\ml{@}};
        \node[lbl]   (succ2) at (2.6,-2.7)  {\ml{\$succ}};
        \node[lbl]   (zero)  at (3.5,-2.7)  {\ml{0}};

        % Tree edges
        \draw (cons2) -- (apa);
        \draw (apa) -- (from);
        \draw (apa) -- (apb);
        \draw (apb) -- (succ1);
        \draw (apb) -- (apc);
        \draw (apc) -- (succ2);
        \draw (apc) -- (zero);

        % First pointer: straight into the CONS root
        \draw[ptr] (s1) -- (1.55,0);

        % Second pointer: from the CONS root, down and around into @c
        \draw[ptr] (cons2.west) to[out=-140, in=180] (apc.west);

        % Third pointer: from the second cell, a wide sweep into @c from below
        \draw[ptr] (s2) to[out=-80, in=-150] (apc.south west);
    \end{tikzpicture}


     &
     % l

     \begin{tikzpicture}[
         baseline=(current bounding box.north),
         font=\sffamily\footnotesize,
         cell/.style={draw, minimum width=1.3cm, minimum height=0.5cm, inner sep=0pt},
         bullet/.style={circle, fill=black, inner sep=1.5pt},
         ptr/.style={-{Stealth[length=2.5mm,width=1.5mm]}, thick},
         apply/.style={inner sep=1pt},
         lbl/.style={inner sep=1pt}
         ]

         % Stack: two cells, each with a bullet at its center
         \node[cell] (c1) at (0,0)    {};
         \node[bullet] (s1) at (c1.center) {};

         % CONS root: stacked double bullet
         \node[bullet] (cons1) at (1.7,0.1)  {};
         \node[bullet] (cons2) at (1.7,-0.1) {};

         % Tree on a regular grid: right child +0.4 (or +0.45) in x, left leaf -0.5 in x, 0.7 down per level
         \node[apply] (apa)   at (2.3,-0.6)  {\ml{@}};
         \node[lbl]   (from)  at (1.8,-1.3)  {\ml{\$from}};
         \node[apply] (apb)   at (2.7,-1.3)  {\ml{@}};
         \node[lbl]   (succ1) at (2.2,-2.0)  {\ml{\$succ}};
         \node[apply] (apc)   at (3.1,-2.0)  {\ml{@}};
         \node[lbl]   (succ2) at (2.6,-2.7)  {\ml{\$succ}};
         \node[lbl]   (zero)  at (3.5,-2.7)  {\ml{0}};

         % Tree edges
         \draw (cons2) -- (apa);
         \draw (apa) -- (from);
         \draw (apa) -- (apb);
         \draw (apb) -- (succ1);
         \draw (apb) -- (apc);
         \draw (apc) -- (succ2);
         \draw (apc) -- (zero);

         % First pointer: straight into the CONS root
         \draw[ptr] (s1) -- (1.55,0);

         % Second pointer: from the CONS root, down and around into @c
         \draw[ptr] (cons2.west) to[out=-140, in=180] (apc.west);
     \end{tikzpicture}

     \\[1cm]
    (j) \ml{CONS} &
    (k) \ml{UPDATE 2} &
    (l) \ml{POP 1} \\
\end{tabular}


%            \end{minipage}%
        }%
    }%

    \caption{\textsf Evaluation of \ml{(\$from (\$succ 0))}}
\end{figure}


Control has now passed to \ml{\$from}. A G-machine compiler will produce the
following G-code for \ml{\$from}, which is executed in sequence:

\simpletitledbox{
    G-code for function \ml{\$from}
}{
    \begin{tabular}{@{}l@{\hspace{2cm}}l@{}}
        \ml{PUSH 0;}            & Push \ml{n} \\
        \ml{PUSHGLOBAL \$succ;} & Push function \ml{\$succ} \\
        \ml{MKAP;}              & Construct \ml{(\$succ n)} \\
        \ml{PUSHGLOBAL \$from;} & Push function \ml{\$from} \\
        \ml{MKAP;}              & Construct \ml{(\$from (\$succ n))} \\
        \ml{PUSH 1;}            & Push \ml{n} \\
        \ml{CONS;}              & Construct \ml{(n : (\$from (\$succ n)))} \\
        \ml{UPDATE 2;}          & Update the root of the redex \\
        \ml{POP 1;}             & Pop the parameter \ml{n} \\
        \ml{UNWIND;}            & Initiate next reduction \\
    \end{tabular}
}

The execution of \ml{\$from} is shown step by step in Figure~18.2. We can make
several observations by examining the code given above:

\begin{numbered}
    \item At the point of entry, the parameter \ml{n} is on top of the stack, and a pointer
    to the root of the redex is immediately below it (Figure~18.2(c)).
    \item Items which are not on top of the stack are addressed \textit{relative to the top of
        the stack}, with the top element having offset zero. For example, the
    \ml{PUSH 1} instruction takes the element next to top in the stack, and pushes
    it onto the stack. Stack items cannot be addressed relative to the base of
    the stack because a reduction takes place at the tip of the spine, with an
    unknown number of vertebrae above. (An alternative would have been
    to assume a frame pointer, and relegate offset calculation to code
    generation time.)
    \item Some instructions take their operands from the stack and put their result
    on the stack in the manner of a zero address machine. \ml{MKAP} and \ml{CONS}
    are examples of such instructions.
\end{numbered}

\noindent
Apart from the last three instructions, the sequence simply constructs an
instance of the body of \ml{\$from} (see Figure~18.2(l)).

The \ml{UPDATE 2} instruction updates the root of the redex with a copy of the
root of the result (there is a slight inefficiency here, since the root of the result
is discarded almost immediately it is constructed; we will address this
efficiency question later). Notice that the G-machine updates the root of the
redex using copying, rather than using indirection nodes (but this is not an
inherent property of the G-machine -- see Section~19.4.4).

The \ml{POP 1} instruction removes the parameters (only one in this case) from
the stack, leaving a pointer to the reduced graph on top of the stack. Finally
\ml{UNWIND} examines the tag of the root node of the reduced graph. In this case it
is a \ml{CONS} cell, so evaluation is complete.

This concludes our example, for now. (Note: in order to reduce the number
of execution steps, the example contains some optimizations which we will
not study until Chapter~20.)

We will now develop the G-machine in a stepwise fashion, beginning with a
very simple implementation, and developing the compilation algorithm and
the G-code together. First, however, we will specify the language from which
we are compiling.

\section{The Source Language for the G-compiler}

The compilation to G-code begins with a program consisting of a number of
supercombinator definitions of the form

\begin{mlcoded}
    \$S x$_1$ x$_2$ \ldots\ x$_n$ $=$ E
\end{mlcoded}

\noindent
where \ml{E} is an expression containing no lambdas, but which may contain \ml{let}s
and \ml{letrec}s. Figure~18.3 gives a reminder of the syntax of expressions. Notice

\boxedfigure{
    \begin{tabular}{@{}lll@{}}
        \ml{<E>} & \typedecl{} & \ml{<constant>} \\
        & \ml{|} & \ml{<identifier>} \\
        & \ml{|} & \ml{<E> <E>} \\
        & \ml{|} & \ml{let <identifier> $=$ <E> in <E>} \\
        & \ml{|} & \ml{letrec <identifier> $=$ <E>} \\
        & & \quad\quad\quad\ml{\ldots} \\
        & & \quad\quad\quad\ml{<identifier> $=$ <E>} \\
        & & \ml{in <E>} \\
    \end{tabular}
}{BNF for syntax of expressions}

\noindent
that the left-hand side of a definition in a \ml{let} or \ml{letrec} can consist only of a
single variable; local function definitions have been removed by lambda-lifting. For example,

\begin{mlcoded}
    let f x $=$ $+$ x 1 \\
    in E
\end{mlcoded}

\noindent
cannot occur. Notice also that we allow only one definition in a \ml{let}. Multiple
definitions can be handled by nested \ml{let}s, and the restriction slightly simplifies
the compiler.

It is worth having a formal description of the syntax, because our compiler
will need to contain a case for each construct. Referring to the syntax enables
us to confirm that all cases have been covered.

To save repetitive work in this chapter we will use a stripped-down set of
built-in functions and constants, shown in Figure~18.4. The stripped-down set
has been chosen to illustrate all the features of the compiler. The operators in
the right-hand column behave exactly like those in the left-hand column.

Assuming that we implement lists with structure tag 1 for \ml{NIL} and 2 for